\chapterhead{5}{Ramps and Effort}{A longer, gentler push for a smaller effort}\label{ch5}

Lifting a heavy box straight up onto a truck is hard. Sliding it up a long plank is easier. The plank is a \textbf{ramp}. A ramp does not make the box lighter. It lets you push a little at a time over a longer path.

\begin{center}
\begin{tikzpicture}
  \draw[line width=1.4pt,color=ink] (0,0) -- (6,0) -- (6,1.6) -- cycle;
  \fill[sky!50,draw=ink,rotate=15.0] (2.2,0.02) rectangle (3.0,0.7);
  \draw[->,line width=1.6pt,color=tomato] (3.0,1.1) -- (3.0,0.2) node[below,font=\sffamily\scriptsize] {gravity};
  \draw[->,line width=1.6pt,color=leaf!70!black] (2.55,0.9) -- (2.35,1.8) node[above,font=\sffamily\scriptsize] {ramp push};
  \draw[->,line width=1.6pt,color=sun!80!black] (3.4,1.1) -- (2.8,0.9) node[left,font=\sffamily\scriptsize] {};
\end{tikzpicture}
\end{center}

\begin{tryit}{How far does the car roll?}
\textbf{You need:} a board, three books, a toy car, a ruler or tape, a pencil.
\begin{enumerate}[leftmargin=1.6em]
\item Lean the board on one book. Let the car go from the top mark without pushing. Measure how far from the bottom of the ramp it stops. Do it three times.
\item Repeat with two books, then three books. Same start mark, same car, same floor.
\item Average each set.
\end{enumerate}
\par\medskip
\renewcommand{\arraystretch}{1.6}
\begin{tabular}{|c|c|c|c|c|}
\hline
\textbf{Books} & \textbf{Roll 1 (cm)} & \textbf{Roll 2 (cm)} & \textbf{Roll 3 (cm)} & \textbf{Average (cm)}\\ \hline
1 & & & & \\ \hline
2 & & & & \\ \hline
3 & & & & \\ \hline
\end{tabular}
\end{tryit}

\begin{worked}{Practice numbers, not from a real test}
One book: 62, 60, 64. Total 186, average $186\div3=62$~cm, spread 4~cm.\\
Two books: 98, 102, 100. Total 300, average 100~cm, spread 4~cm.\\
Three books: 130, 132, 134. Total 396, average 132~cm, spread 4~cm.\\
A steeper ramp sent the car farther each time. The gaps between averages (38 and 32~cm) are much bigger than the spread (4~cm), so the pattern is not just wobble.
\end{worked}

\begin{multix}[hmm]\multixbot[hmm]\tcblower
The car went farther on the steeper ramp, so a steeper ramp means a bigger force of gravity on the car. Gravity got stronger! Did it?
\end{multix}

\trust{We changed only the number of books, kept the start mark, the car and the floor the same, and rolled three times for each ramp.}
{The averages rose with each book and the gaps were much larger than the spread of repeats.}
{Whether it holds for a heavier car, a rough floor or a much steeper ramp. We tried three heights.}

\begin{keepbox}
We drew the car on the ramp as a few arrows. \textbf{We kept} \wline[3.4cm] \textbf{and we set aside} \wline[3.4cm].
\end{keepbox}

\clearpage

\subhead{Pull to lift, or pull up a ramp}
Hang a small bag of coins from a rubber band. Lift it straight up and see how far the band stretches. Then pull the same bag up a ramp with the band. The stretch is a stand-in for how hard you pull.

\begin{tryit}{Straight up or up the ramp?}
\textbf{You need:} a rubber band, a bag of coins, a board, books, a ruler.
\begin{enumerate}[leftmargin=1.6em]
\item Lift the bag slowly straight up. Measure the stretch while it moves.
\item Lean the board to the same height. Pull the bag slowly up the board. Measure the stretch.
\item Use a longer board to reach the same height. Measure again.
\end{enumerate}
\par\medskip
\renewcommand{\arraystretch}{1.6}
\begin{tabular}{|l|c|}
\hline
\textbf{Way up} & \textbf{Stretch (cm)}\\ \hline
Straight up & \\ \hline
Short ramp & \\ \hline
Long ramp & \\ \hline
\end{tabular}
\end{tryit}

\begin{worked}{Practice numbers}
Straight lift: 6.0~cm stretch. A 40~cm ramp: 3.8~cm. An 80~cm ramp: 2.5~cm.\\
Longer ramp, smaller pull. But the bag travels 80~cm instead of 40~cm, so the pull lasts twice as far.
\end{worked}

\begin{center}
\begin{tikzpicture}
  \node[bigcir,minimum size=2.3cm] (a) at (0,0) {}; \node[font=\sffamily\scriptsize,text width=1.9cm,align=center] at (0,0) {steep\\short ramp\\pull 3.8};
  \node[bigcir,minimum size=2.3cm] (b) at (6,0) {}; \node[font=\sffamily\scriptsize,text width=1.9cm,align=center] at (6,0) {gentle\\long ramp\\pull 2.5};
  \draw[popA] (a) -- node[above,note] {make it longer} (b);
  \node[note] at (3,-1.5) {\textbf{kept:} the direction of each pull. \textbf{lost:} the real shape of the board.};
\end{tikzpicture}
\end{center}

\trust{We pulled the same bag the same height three ways and read the band each time.}
{The stretch got smaller as the ramp got longer: 6.0, 3.8, 2.5.}
{That the pattern is exact. Friction on the board also changes the pull, and we measured once for each.}

\begin{keepbox}
A ramp trades push for distance. \textbf{We kept} \wline[3.4cm] \textbf{and we set aside} \wline[3.4cm].
\end{keepbox}

\clearpage

\begin{puzzle}{Average the rolls}
A car rolled 55, 57 and 59~cm on a ramp with one book.\\
Total \wline[1.6cm] Average \wline[1.6cm]~cm
\end{puzzle}

\begin{puzzle}{Predict the roll}
On two books the car averaged 100~cm. On three books it averaged 132~cm. Without testing, guess the average for four books.\\
My guess: \wline[2cm]~cm. Why is it only a guess? \wline[3cm]
\end{puzzle}

\begin{puzzle}{Which ramp?}
You must lift a box 1~m onto a truck. One ramp is 2~m long and one is 4~m long. Which needs a smaller push? \wline[2cm]\\
What do you pay for the smaller push? \wline[3cm]
\end{puzzle}

\begin{puzzle}{Unfair ramp test}
Ana rolled a car down a steep ramp on a carpet. Ben rolled the same car down a gentle ramp on a smooth floor. Ben's car went farther. Ana says ``a gentle ramp is better''. What is wrong with the test?\\
\rule{\linewidth}{0.5pt}
\end{puzzle}

\begin{puzzle}{Draw the arrows}
A sled sits still on a snowy slope. Draw three arrows from the dot: gravity, the slope's push, and friction.
\begin{center}
\begin{tikzpicture}
  \draw[line width=1.4pt,color=ink] (0,0) -- (6,0) -- (6,2) -- cycle;
  \fill[ink] (3,1.0) circle (3pt);
\end{tikzpicture}
\end{center}
Is anything missing? \wline[3cm]
\end{puzzle}

\begin{puzzle}{Slow down}
A toy car rolls down a ramp and then across a rug. On the ramp it speeds up. On the rug it slows. Which arrow gets bigger on the rug? \wline[2.5cm]
\end{puzzle}
\clearpage

\begin{puzzle}{Gaps}
Averages for one, two and three books were 62, 100 and 132~cm. Find the gap from one to two books \wline[1.6cm] and from two to three \wline[1.6cm].
\end{puzzle}

\begin{puzzle}{How much easier?}
Straight lift: 6.0. 40~cm ramp: 3.8. 80~cm ramp: 2.5. By how much smaller is the pull on the 40~cm ramp than the lift? \wline[1.6cm]\\
On the 80~cm ramp? \wline[1.6cm]
\end{puzzle}

\begin{puzzle}{Speeding up}
A toy car rolls down a steep ramp, then a gentle ramp. On which does it speed up more? \wline[3cm]\\
Which arrow along the slope is bigger on the steep one? \wline[3cm]
\end{puzzle}

\begin{puzzle}{Ramp hunt}
Find three ramps near you. Write where each is and who it is for.\\
\rule{\linewidth}{0.5pt}\\[10pt]\rule{\linewidth}{0.5pt}\\[10pt]\rule{\linewidth}{0.5pt}
\end{puzzle}
\clearpage
